%!TEX root=../../note
\subsection{Proving Universally Quantified Statements}
A \textbf{proposition} is a mathematical statement established by proof. A \textbf{theorem}
is a particularly significant proposition. A \textbf{lemma} is a subsidiary proposition.
A \textbf{corollary} is a proposition that follows almost immediately from a theorem.

\begin{example}
    The law of cosines is a theorem: in a triangle with sides $a, b, c$, where $\theta$ is the angle between
    $a$ and $b$,
    \begin{equation*}
        c^2 = a^2 + b^2 - 2ab\cos \theta.
    \end{equation*}
    The Pythagorean theorem is a corollary: taking $\theta = \pi/2$ gives $\cos\theta = 0$, so $c^2 = a^2 + b^2$.
\end{example}

\subsubsection{An Arbitrary Element}
To prove $\forall x \in S, P(x)$, let $x$ be an \emph{arbitrary} element of $S$, using only the fact that
$x \in S$, and show that $P(x)$ is true. Since nothing special about $x$ was assumed, the argument
works for every element of $S$.

The following inequalities are proved by rewriting a difference as a square.
\begin{proposition}
    \begin{equation*}
        \forall x \in \R, x^2 + 1 \geq 2x.
    \end{equation*}
\end{proposition}
\begin{remark}[Proof idea]
    Moving everything to one side gives $x^2 - 2x + 1 \geq 0$, and the left side is the perfect square
    $(x - 1)^2$. Since squares of real numbers are never negative, we start from $(x-1)^2 \geq 0$
    and work forward.
\end{remark}
\begin{proof}
    Let $x$ be an arbitrary real number. Therefore, $x - 1$ must also be a real number, and hence
    \begin{equation*}
        (x - 1)^2 \geq 0.
    \end{equation*}

    Expanding the left side, we get
    \begin{equation*}
        x^2 - 2x + 1 \geq 0.
    \end{equation*}

    Adding $2x$ to both sides yields
    \begin{equation*}
        x^2 + 1 \geq 2x.
    \end{equation*}
\end{proof}

\begin{proposition}
    \begin{equation*}
        \forall x, y \in \R, x^4 + x^2y + y^2 \geq 5x^2y - 3y^2.
    \end{equation*}
\end{proposition}
\begin{remark}[Proof idea]
    Moving everything to the left side gives $x^4 - 4x^2y + 4y^2 \geq 0$, and
    $x^4 - 4x^2y + 4y^2 = (x^2 - 2y)^2$. This explains the choice of $x^2 - 2y$ below.
\end{remark}
\mdfsetup{nobreak=true}
\begin{proof}
    Let $x$ and $y$ be arbitrary real numbers. Then $x^2 - 2y$ is also a real number. So
    \begin{equation*}
        (x^2 - 2y)^2 \geq 0.
    \end{equation*}

    Expanding, we get
    \begin{equation*}
        x^4 - 4x^2y + 4y^2 \geq 0.
    \end{equation*}

    Adding $5x^2y$ to both sides, we have
    \begin{equation*}
        x^4 + x^2y + 4y^2 \geq 5x^2y.
    \end{equation*}

    Now, subtract $3y^2$ from both sides:
    \begin{equation*}
        x^4 + x^2y + y^2 \geq 5x^2y - 3y^2.
    \end{equation*}
\end{proof}
\mdfsetup{nobreak=false}

\subsubsection{Proof by Cases}
\begin{proposition}
    \begin{equation*}
        \forall x, y \in \R, \max(x, y) = \frac{x + y + \abs{x - y }}{2}.
    \end{equation*}
\end{proposition}
\begin{remark}[Proof idea]
    Both $\max(x, y)$ and $\abs{x - y}$ are defined by cases depending on which of $x, y$ is larger,
    so we split into the cases $x \geq y$ and $x < y$. These cases cover every pair of real numbers.
\end{remark}
\begin{proof}
    Let $x$ and $y$ be arbitrary real numbers.

    \textbf{Case 1: $x\geq y$.} Then $\max(x, y) = x$ and $x - y \geq 0$, so $\abs{x - y } = x - y$.

    Then
    \begin{equation*}
        \frac{x + y + \abs{x - y }}{2} = \frac{x + y + x - y }{2 } = \frac{2x }{2 } = x = \max(x, y).
    \end{equation*}

    So the claim holds when $x \geq y$.

    \textbf{Case 2: $x<y$.} Then $\max(x, y) = y$ and $x - y < 0$, so
    \begin{equation*}
        \abs{x - y } = -(x - y) = y - x.
    \end{equation*}

    Then
    \begin{equation*}
        \frac{x + y + \abs{x - y }}{2} = \frac{x + y + y - x }{2} = \frac{2y}{2} = y = \max(x, y).
    \end{equation*}

    The claim holds in both cases, so it holds for all $x, y \in \R$.
\end{proof}

\subsubsection{Disproving a Universal Statement}
Since $\neg \paren{\forall x \in S, P(x)} \equiv \exists x \in S, \neg P(x)$, to disprove a universal
statement it suffices to give one \textbf{counterexample}: a specific $x \in S$ for which $P(x)$ is false.

\begin{example}
    Disprove:
    \begin{equation*}
        \forall x \in \R, (x^2 - 1)^2 > 0.
    \end{equation*}
\end{example}
\noindent\emph{Solution.} Take $x = 1$. Then
\begin{equation*}
    \paren{x^2 - 1}^2 = (1 - 1)^2 = 0 \not> 0.
\end{equation*}
So $x = 1$ is a counterexample, and the statement is false.

\subsubsection*{In practice}
\begin{itemize}
    \item \textbf{Universal proof.} Start with an arbitrary element of the
    stated domain. Use the hypotheses to derive the claim without choosing
    a special value of that element.
    \item \textbf{Cases.} Choose cases that cover the entire domain and
    prove the same conclusion in each. For absolute values, the sign of the
    expression inside determines the cases.
    \item \textbf{Counterexample.} Give one element in the domain and
    verify that it fails the claimed condition.
\end{itemize}

\subsubsection*{Exercises}
\begin{enumerate}
    \item Prove $\forall \theta \in \R, \sin 3\theta = 3\sin \theta - 4\sin^3 \theta$.
    \item Prove $\forall x \in \R, \abs{x - 3 } + 2\abs{x + 2} \geq 5$.
\end{enumerate}
